\section{Semantic Interoperability}
\label{sec:semantic}

Semantic interoperability is the single rubric axis on which no
surveyed protocol \emph{ships} a level-3 mechanism
(\S\ref{sec:comparison:matrix}, observation 3). The cluster
sits at level 1, free-text descriptions over typed schemas.
A2A, MCP, ACP, AGNTCY, and Coral all occupy this anchor. ANP at
level 2 (JSON-LD with the Agent Description Protocol) is the only
protocol that ships an ontology-base mechanism. LOKA at level 3
is explicitly aspirational; the Polyglot Intent Engine and
capability ontologies that the LOKA white paper~\cite{spec_loka}
describes are not implemented at the cutoff. This gap is
structural in the precise sense that no contemporary protocol's
roadmap promises to close it within the next revision cycle.
The field is unsure that closing it is the right objective.

This section unpacks the gap. Section~\ref{sec:semantic:syntactic}
explains why level 1 is the local equilibrium for LLM-era
protocols even though it was a known dead-end in the FIPA-ACL era.
Section~\ref{sec:semantic:capabilities} reviews the
capability-description schemas that the contemporary protocols
do ship, and the typed-schema-plus-free-text-description pattern
they share. Section~\ref{sec:semantic:ontology} treats the
ontology and skill-alignment proposals in the
literature~\cite{semantic_driven, beyond_message_passing} and
locates LOKA's design within them. Section~\ref{sec:semantic:fragmentation}
closes with the schema-fragmentation risk that the cluster's
free-text equilibrium is silently producing.

\subsection{Beyond syntactic message passing}
\label{sec:semantic:syntactic}

The contemporary protocols specify message \emph{syntax} with
considerable rigour. JSON-RPC 2.0 envelopes, typed Part vocabularies,
JSON Schema for tool input, content-type negotiation: these
are well-specified and uniformly applied. What none of the
industry-led protocols specifies, beyond a brief free-text
description, is what the message \emph{means}: what the
\texttt{find\_billing\_record} skill returns, what units the
\texttt{amount} field is in, what the agent should do if the
input parameter falls outside an implicit acceptable range.
Meaning is left to the natural-language description and to
the LLM-receiver's ability to interpret it.

The argument for this equilibrium (which we developed in
\S\ref{sec:background:llm-disruption}) is that LLM agents can
recover from semantic mismatches through inference and clarifying
questions in ways that FIPA-era deterministic parsers could not.
Two LLM agents with reasonable defaults can interoperate over a
thin envelope without explicit ontology alignment, and the
flexibility this affords is a real and significant advantage.
Agent designers can author skills without committing to a shared
vocabulary.

The argument against the equilibrium is that ``recovery through
inference'' is exactly the surface that prompt injection
(\S\ref{sec:security:injection}), capability escalation
(\S\ref{sec:security:authz}), and contextual-integrity violations
(\S\ref{sec:security:privacy}) attack. A free-text skill
description that an LLM caller interprets ``correctly enough''
under normal conditions can be re-interpreted under
adversarial conditions to authorise capabilities the
description did not intend. The \emph{Beyond Message Passing}
analysis~\cite{beyond_message_passing} treats this argument
in depth and concludes that the free-text equilibrium is
\emph{fundamentally insufficient} for
cross-organizational interoperability under any threat model
the contemporary literature takes seriously. The
\emph{Semantic-Driven AI Agent Communications} line~\cite{semantic_driven}
reaches a similar conclusion from a different starting point.
It argues that the inference cost of recovering meaning through
LLM reasoning at every message boundary is operationally
unsustainable at scale, regardless of the security argument.

The protocol-design implication is that level 1 is not a
sustainable equilibrium even on its own terms. The question
the field has not answered is what level 2 should look like
that captures enough meaning to address the threats and the
inference cost without re-introducing FIPA-SL's authoring
overhead.

\subsection{Capability description schemas}
\label{sec:semantic:capabilities}

The capability-description schemas the contemporary protocols
ship cluster around four design choices that are worth
enumerating because the cluster's apparent convergence on a
common pattern hides important differences.

\paragraph{A2A's typed Skill objects.}
A2A's AgentCard exposes a \texttt{skills} array whose elements
are typed Skill objects with \texttt{id},
\texttt{name}, \texttt{description}, \texttt{tags}, and
optional \texttt{inputSchema} and \texttt{outputSchema}
fields~\cite{spec_a2a}. The \texttt{description} field is
free text; the schemas, when present, are JSON Schema. The
de-facto convention in deployments is to omit the schemas
and rely on the description, which is the failure mode that
prompts the rubric's level-1 score. The A2A working group has
discussed elevating JSON Schema from optional to
\emph{recommended} for the next revision, which would not
move the rubric score but would meaningfully reduce the
free-text surface in production deployments.

\paragraph{MCP's tool input contracts.}
MCP's \texttt{tools} list exposes each tool with
\texttt{name}, \texttt{description}, and a required
\texttt{inputSchema}~\cite{spec_mcp}. The schema is JSON
Schema, the description is free text, and the implicit
contract is the model that A2A, ACP, and AGNTCY all follow.
The schema constrains the inputs the tool will accept. The
description tells the LLM what the tool does.
MCP's level 2 on Axis 7 (multi-modal handling) and level 1
on Axis 9 reflect the same architectural choice: typed
inputs (Axis 7 typing for content blocks; Axis 9 typing for
parameters) without typed semantics.

\paragraph{ACP's multipart capability bundles.}
ACP's capability-description schema differs from MCP and A2A
in carrying capability descriptions as multipart MIME
bundles rather than JSON Schema~\cite{spec_acp}. The
operational advantage is that the schema can carry richer
content (example inputs, sample outputs, supplementary
documentation) without the brittleness of embedding it all
in a JSON string. The semantic advantage is small: the
multipart bundle is still authored in free text. The
rubric score is the same as A2A's and MCP's.

\paragraph{AGNTCY's Open Agentic Schema Format (OASF).}
The OASF specification~\cite{spec_agntcy} is the most
explicit capability-descriptor effort of the cluster. OASF
defines a schema for agent descriptors that is
intentionally compatible with both A2A AgentCards and MCP
server manifests, with the design intent of allowing a
single directory to serve heterogeneous agents. OASF
descriptors include capability descriptions, but the
descriptions are still free text within typed fields.
OASF's contribution is at Axis 1 (federated directory
with structured descriptors) rather than at Axis 9; the
schema unifies the meta-structure of descriptors across
protocols, not the semantics of capability descriptions.

\paragraph{The shared pattern.}
All four protocols share the pattern: \emph{typed envelope,
typed parameters, free-text semantics}. The pattern is
deliberate. Each protocol's designers considered
alternative patterns and chose this one. The
convergence is evidence that it is the local equilibrium of
the contemporary design space. The question for level 2 is
what changes the equilibrium.

\subsection{Ontology and skill alignment}
\label{sec:semantic:ontology}

The proposals in the literature for moving beyond level 1
divide into two families.

\paragraph{Ontology-base proposals.}
The first family attempts to anchor capability descriptions
in a shared ontology of concepts. ANP's JSON-LD
serialization~\cite{spec_anp} is the canonical contemporary
example, scoring level 2 on Axis 9. The ANP AgentDescription
document uses JSON-LD's
\texttt{@context} mechanism~\cite{w3c_json_ld} to bind
capability names to URIs in a shared vocabulary; agents
that understand the
vocabulary can match capabilities by URI rather than by
free-text similarity. The cost is authoring overhead: every
capability must be located in the vocabulary, and the
vocabulary itself must be maintained against drift.

That cost is not hypothetical, and the evidence arrived
from the one protocol that had paid it. Post-cutoff, ANP's
1.1 release line moved the Agent Description document away
from JSON-LD to plain JSON, while the Discovery
specification kept JSON-LD and its \texttt{@context}
mechanism~\cite{spec_anp_v11}. The stated reason was
readability, and the community describes the Linked Data
philosophy as retained in the design even where the
serialization is not. The survey does not restate the Axis
9 score, which is fixed at the 2026-Q1 cutoff. The
direction still matters for this section's argument. The
only shipping level-2 mechanism in the surveyed set
narrowed its ontology commitment at the description layer,
which is the layer a capability consumer reads. So the
authoring overhead above is not merely a cost a designer
weighs in the abstract. It is a cost that a protocol
carrying the mechanism in production chose to stop paying
in the document where it was most visible. Any proposal to
move the field from level 1 to level 2 must answer that,
and a proposal that assumes ontology adoption is
monotonic is contradicted by the one case study
available.

The LOKA design takes the ontology proposal furthest among
the surveyed protocols. The LOKA capability ontologies
(\S 3.1) commit to a typed
capability hierarchy with intent schemas; the Polyglot
Intent Engine (\S 4.5.2) is the proposed mechanism for
translating between the LOKA ontology and FIPA-ACL, A2A
AgentCards, and the Open Voice Network's intent
specifications. The mechanism is aspirational and no
implementation exists at the cutoff. The design is still
the most explicit articulation in the surveyed set of what
a level-3 ontology-anchored protocol would look like.

The \emph{Beyond Message Passing} proposal~\cite{beyond_message_passing}
generalises the ontology-base direction into a
vocabulary-and-matcher architecture: protocols commit to a
shared vocabulary of capabilities, intents, and outcomes;
agents that author skills locate them in the vocabulary;
matchers at the protocol layer route requests to capabilities
by ontology proximity rather than by free-text similarity.
The proposal is at the pre-specification stage; it has not
been adopted into any protocol's roadmap.

\paragraph{Semantic-matchmaker proposals.}
The second family is the semantic-matchmaker pattern. Rather
than committing to a shared ontology authored in advance, the
matchmaker proposes that agents author capabilities in free
text and that a semantic-similarity matcher (typically
embedding-based) at the discovery layer maps requests to
capabilities by embedding distance. The
\emph{Semantic-Driven AI Agent Communications} analysis
\cite{semantic_driven} surveys the matchmaker space and
identifies its advantages: no ontology to maintain, no
authoring overhead, and graceful degradation when the
request doesn't match any single capability cleanly.

The cost is what the security analysis
(\S\ref{sec:security:injection}) already identifies:
embedding-based matching is exactly the surface that
adversarial natural-language content attacks. A
capability description that ranks highly under semantic
similarity to a malicious request is invoked by the
matchmaker without further constraint; the rubric's level
2 anchor (\emph{ontology-base or semantic-discovery indexing})
intentionally groups both approaches because both fail in
the same way under adversarial conditions.

\paragraph{The FIPA-SL retrospective.}
A useful question to ask of both families is why FIPA-SL's
full ontology commitment did not survive the early 2000s.
The historical answer is twofold: the authoring cost was
high (ontologies had to be defined, then maintained, then
agreed across deployments) and the matching infrastructure
was brittle (translation services between ontologies were
the bottleneck of every cross-organizational deployment).
Both arguments still apply, modulated by the LLM era's
capacity to bridge ontology mismatches through inference.
The contemporary proposals take this lesson seriously: ANP's
JSON-LD vocabulary is deliberately lighter than FIPA-SL;
LOKA's Polyglot Intent Engine is a translator-of-translators
that takes ontology mismatch as a first-class operational
concern.

\subsection{The schema fragmentation risk}
\label{sec:semantic:fragmentation}

The free-text equilibrium produces a quieter risk that the
threat-coverage literature does not catalogue but that
practitioners report at scale: \emph{schema fragmentation}.
Two organisations deploying A2A agents with
free-text skill descriptions will, with high probability,
adopt slightly different conventions for what a billing
query looks like, what a customer-record schema includes,
what units \texttt{amount} is denominated in. Each
organisation's deployments work internally; cross-organizational
calls fail in ways that the protocol cannot diagnose. The
LLM-receiver's interpretation is the only place where the
mismatch surfaces, and the failure mode is silent
mis-execution rather than explicit error.

Three consequences are visible in the
deployment literature at the cutoff.

\paragraph{The inter-organisational call cost.}
Cross-organisational A2A calls have been documented in the
AAIF interoperability working group's public
discussions~\cite{lf_aaif} to succeed at substantially
lower rates than intra-organisational calls, with the gap
attributed overwhelmingly to schema mismatch rather than to
transport failure or authentication failure. The mismatch is
between skills with the same name and similar descriptions
that interpret parameters differently. The empirical finding
is consistent with the AgentLeak benchmark's identification
of inter-agent communication as the dominant
sensitive-content failure surface
(\S\ref{sec:security:privacy})~\cite{agentleak}: the same
parameter-interpretation disagreement that causes a
billing query to fail across organisations is what makes a
clarifying-question retry leak content the principal did
not authorise the remote agent to receive.

\paragraph{The retrofit cost.}
Organisations that begin with free-text skill descriptions
and later attempt to migrate to JSON Schema or OASF
descriptors report substantial retrofit cost. The reason is
that production deployments grow callers that depend on
the free-text behaviour, and adding a schema after the
fact requires either matching the schema to the de-facto
behaviour (which removes the schema's value as a contract)
or breaking the callers (which removes the value of
migration). The retrofit pattern is sometimes referred to
in the deployment literature as ``the second-system schema
penalty.''

\paragraph{The matchmaker drift.}
Deployments that rely on embedding-based semantic
matchmakers report drift: the same query that returned
capability X under last quarter's embedding model returns
capability Y under this quarter's, because the embedding
space has shifted with the underlying foundation model.
The cost is that capability invocation is not
reproducible across foundation-model upgrades; the
audit-log trail must record the embedding-model
identifier to be useful for after-the-fact diagnostics.
No contemporary protocol's audit specification
(\S\ref{sec:security:tee}) carries this discipline.

\paragraph{The protocol-versus-orchestrator question.}
The dispatch question that closes this section is whether
the schema-fragmentation risk should be addressed at the
protocol layer or at the orchestrator layer. The
\emph{Beyond Message Passing} position is that it must
be at the protocol layer because cross-organisational
deployments cannot assume a common orchestrator. The
\emph{Semantic-Driven} position is that the orchestrator
is the natural place for matchmaker logic because the
matcher needs visibility into the principal's intent.
The survey takes no view, since the question is properly an
open problem (\S\ref{sec:openproblems}). We observe
that the matchmaker drift consequence above is
operationally visible to whichever layer owns the
matcher, and that the audit-log discipline necessary to
diagnose it is a protocol concern regardless of where the
matcher itself lives.

The principal conclusion of \S\ref{sec:semantic} is that
the contemporary protocols' level-1 cluster is a known
local equilibrium with known costs that the field has not
yet decided how to escape. The escape route is the
single largest research-and-standardization question on
the survey's open-problems agenda
(\S\ref{sec:openproblems}); the candidate answers
catalogued here will reappear there as the basis for the
research-and-standardization recommendations.
